% Options for packages loaded elsewhere
\PassOptionsToPackage{unicode}{hyperref}
\PassOptionsToPackage{hyphens}{url}
\PassOptionsToPackage{dvipsnames,svgnames,x11names}{xcolor}
%
\documentclass[
  10pt,
  a4paper,
]{article}
\usepackage{amsmath,amssymb}
\usepackage{iftex}
\ifPDFTeX
  \usepackage[T1]{fontenc}
  \usepackage[utf8]{inputenc}
  \usepackage{textcomp} % provide euro and other symbols
\else % if luatex or xetex
  \usepackage{unicode-math} % this also loads fontspec
  \defaultfontfeatures{Scale=MatchLowercase}
  \defaultfontfeatures[\rmfamily]{Ligatures=TeX,Scale=1}
\fi
\usepackage{lmodern}
\ifPDFTeX\else
  % xetex/luatex font selection
\fi
% Use upquote if available, for straight quotes in verbatim environments
\IfFileExists{upquote.sty}{\usepackage{upquote}}{}
\IfFileExists{microtype.sty}{% use microtype if available
  \usepackage[]{microtype}
  \UseMicrotypeSet[protrusion]{basicmath} % disable protrusion for tt fonts
}{}
\makeatletter
\@ifundefined{KOMAClassName}{% if non-KOMA class
  \IfFileExists{parskip.sty}{%
    \usepackage{parskip}
  }{% else
    \setlength{\parindent}{0pt}
    \setlength{\parskip}{6pt plus 2pt minus 1pt}}
}{% if KOMA class
  \KOMAoptions{parskip=half}}
\makeatother
\usepackage{xcolor}
\usepackage[margin=2.1cm]{geometry}
\usepackage{color}
\usepackage{fancyvrb}
\newcommand{\VerbBar}{|}
\newcommand{\VERB}{\Verb[commandchars=\\\{\}]}
\DefineVerbatimEnvironment{Highlighting}{Verbatim}{commandchars=\\\{\}}
% Add ',fontsize=\small' for more characters per line
\newenvironment{Shaded}{}{}
\newcommand{\AlertTok}[1]{\textcolor[rgb]{1.00,0.00,0.00}{\textbf{#1}}}
\newcommand{\AnnotationTok}[1]{\textcolor[rgb]{0.38,0.63,0.69}{\textbf{\textit{#1}}}}
\newcommand{\AttributeTok}[1]{\textcolor[rgb]{0.49,0.56,0.16}{#1}}
\newcommand{\BaseNTok}[1]{\textcolor[rgb]{0.25,0.63,0.44}{#1}}
\newcommand{\BuiltInTok}[1]{\textcolor[rgb]{0.00,0.50,0.00}{#1}}
\newcommand{\CharTok}[1]{\textcolor[rgb]{0.25,0.44,0.63}{#1}}
\newcommand{\CommentTok}[1]{\textcolor[rgb]{0.38,0.63,0.69}{\textit{#1}}}
\newcommand{\CommentVarTok}[1]{\textcolor[rgb]{0.38,0.63,0.69}{\textbf{\textit{#1}}}}
\newcommand{\ConstantTok}[1]{\textcolor[rgb]{0.53,0.00,0.00}{#1}}
\newcommand{\ControlFlowTok}[1]{\textcolor[rgb]{0.00,0.44,0.13}{\textbf{#1}}}
\newcommand{\DataTypeTok}[1]{\textcolor[rgb]{0.56,0.13,0.00}{#1}}
\newcommand{\DecValTok}[1]{\textcolor[rgb]{0.25,0.63,0.44}{#1}}
\newcommand{\DocumentationTok}[1]{\textcolor[rgb]{0.73,0.13,0.13}{\textit{#1}}}
\newcommand{\ErrorTok}[1]{\textcolor[rgb]{1.00,0.00,0.00}{\textbf{#1}}}
\newcommand{\ExtensionTok}[1]{#1}
\newcommand{\FloatTok}[1]{\textcolor[rgb]{0.25,0.63,0.44}{#1}}
\newcommand{\FunctionTok}[1]{\textcolor[rgb]{0.02,0.16,0.49}{#1}}
\newcommand{\ImportTok}[1]{\textcolor[rgb]{0.00,0.50,0.00}{\textbf{#1}}}
\newcommand{\InformationTok}[1]{\textcolor[rgb]{0.38,0.63,0.69}{\textbf{\textit{#1}}}}
\newcommand{\KeywordTok}[1]{\textcolor[rgb]{0.00,0.44,0.13}{\textbf{#1}}}
\newcommand{\NormalTok}[1]{#1}
\newcommand{\OperatorTok}[1]{\textcolor[rgb]{0.40,0.40,0.40}{#1}}
\newcommand{\OtherTok}[1]{\textcolor[rgb]{0.00,0.44,0.13}{#1}}
\newcommand{\PreprocessorTok}[1]{\textcolor[rgb]{0.74,0.48,0.00}{#1}}
\newcommand{\RegionMarkerTok}[1]{#1}
\newcommand{\SpecialCharTok}[1]{\textcolor[rgb]{0.25,0.44,0.63}{#1}}
\newcommand{\SpecialStringTok}[1]{\textcolor[rgb]{0.73,0.40,0.53}{#1}}
\newcommand{\StringTok}[1]{\textcolor[rgb]{0.25,0.44,0.63}{#1}}
\newcommand{\VariableTok}[1]{\textcolor[rgb]{0.10,0.09,0.49}{#1}}
\newcommand{\VerbatimStringTok}[1]{\textcolor[rgb]{0.25,0.44,0.63}{#1}}
\newcommand{\WarningTok}[1]{\textcolor[rgb]{0.38,0.63,0.69}{\textbf{\textit{#1}}}}
\setlength{\emergencystretch}{3em} % prevent overfull lines
\providecommand{\tightlist}{%
  \setlength{\itemsep}{0pt}\setlength{\parskip}{0pt}}
\setcounter{secnumdepth}{5}
\ifLuaTeX
  \usepackage{selnolig}  % disable illegal ligatures
\fi
\usepackage{bookmark}
\IfFileExists{xurl.sty}{\usepackage{xurl}}{} % add URL line breaks if available
\urlstyle{same}
\hypersetup{
  colorlinks=true,
  linkcolor={Maroon},
  filecolor={Maroon},
  citecolor={Blue},
  urlcolor={Blue},
  pdfcreator={LaTeX via pandoc}}

\author{}
\date{}

\begin{document}

{
\hypersetup{linkcolor=}
\setcounter{tocdepth}{3}
\tableofcontents
}
\section{ThreatForge / DDTA Decision candidates after Revision 7
regression}\label{threatforge-ddta-decision-candidates-after-revision-7-regression}

\begin{quote}
\textbf{Non-canonical research rewrite.}

This document does not patch ThreatForge. It rewrites the sixteen
historical construction slots after the Revision 7 regression. A slot
may remain a ThreatForge Decision, move to another semantic owner/layer,
keep its parent unresolved, or cease to be a standalone Decision.
\end{quote}

\subsection{Reading rule}\label{reading-rule}

The canonical Decision-specific semantic shape remains:

\begin{Shaded}
\begin{Highlighting}[]
\NormalTok{Decision}
\NormalTok{|{-} title}
\NormalTok{|{-} context}
\NormalTok{|{-} decision}
\NormalTok{\textasciigrave{}{-} consequences}
\end{Highlighting}
\end{Shaded}

Each section below also carries \textbf{study annotations}
(\texttt{Disposition}, \texttt{Semantic\ owner},
\texttt{Proposed\ parent}, \texttt{DDTA\ contract\ supported}). These
annotations exist only to evaluate layering and ownership. They are
\textbf{not proposed fields of the Decision model}.

\subsection{Corrected construction
corpus}\label{corrected-construction-corpus}

\subsection{MR-0001 / ADR-0001}\label{mr-0001-adr-0001}

\textbf{Disposition:} Relocated Decision candidate

\textbf{Semantic owner:} DDTA method

\textbf{Proposed parent:} Not assigned inside the ThreatForge MR
hierarchy

\textbf{DDTA contract supported / related:} Document classification
policy (method-level candidate)

\subsubsection{Title}\label{title}

Diátaxis come classificazione comune della funzione documentale

\subsubsection{Context}\label{context}

Una documentazione governata usata in progetti diversi deve permettere a
lettori e strumenti di riconoscere la funzione editoriale di un
documento senza dover apprendere ogni volta una classificazione
proprietaria. Il metodo deve quindi scegliere se proporre una
classificazione comune oppure lasciare interamente ai singoli progetti
la distinzione fra materiale per apprendere, eseguire un compito,
consultare un riferimento o comprendere un argomento.

\subsubsection{Decision}\label{decision}

Il metodo DDTA adotta Diátaxis come classificazione comune della
funzione editoriale dei documenti, distinguendo tutorial, how-to guide,
reference ed explanation secondo il bisogno del lettore. La
classificazione editoriale non determina da sola l'autorità semantica
del contenuto, che rimane definita dal modello governato e dalle fonti
che possiedono l'informazione.

\subsubsection{Consequences}\label{consequences}

\begin{itemize}
\tightlist
\item
  Progetti diversi possono condividere una classificazione leggibile
  senza inventare una tassonomia editoriale locale.
\item
  La classificazione richiede giudizio editoriale e non può essere
  ridotta alla sola posizione del file.
\item
  Il supporto ThreatForge a Diátaxis diventa una questione di
  conformance/tooling distinta dalla scelta di metodo.
\end{itemize}

\subsubsection{Study note}\label{study-note}

Resta una candidata Decision di metodo; la gerarchia DDTA che dovrebbe
possederla non è definita in questo checkpoint.

\begin{center}\rule{0.5\linewidth}{0.5pt}\end{center}

\subsection{MR-0001 / ADR-0002}\label{mr-0001-adr-0002}

\textbf{Disposition:} Relocated Decision candidate

\textbf{Semantic owner:} DDTA method / cross-document governance

\textbf{Proposed parent:} Not assigned inside the ThreatForge MR
hierarchy

\textbf{DDTA contract supported / related:} Shared terminology
governance (status still open)

\subsubsection{Title}\label{title-1}

Vocabolario governato come fonte unica del linguaggio condiviso

\subsubsection{Context}\label{context-1}

Concetti condivisi possono ricorrere in documenti, progetti, sessioni e
strumenti differenti. Se definizioni e sinonimi vengono ricostruiti
localmente, il corpus può divergere semanticamente anche quando ogni
singolo documento rimane formalmente valido. Occorre quindi scegliere se
il significato dei termini condivisi dipenda da convenzioni locali
oppure da conoscenza governata comune.

\subsubsection{Decision}\label{decision-1}

DDTA usa, quando il profilo/metodo adottato lo richiede, una fonte
governata unica per i termini condivisi e per il loro significato. Le
definizioni comuni vengono referenziate o risolte dalla fonte che le
possiede invece di essere copiate cumulativamente nei documenti
discendenti.

\subsubsection{Consequences}\label{consequences-1}

\begin{itemize}
\tightlist
\item
  Persone e strumenti possono ricostruire il significato senza dipendere
  dalla memoria di una sessione precedente.
\item
  Le modifiche terminologiche hanno un proprietario governato invece di
  richiedere copie sincronizzate.
\item
  Resta da stabilire se il controlled vocabulary sia universale,
  opzionale, profile-specific o soltanto raccomandato.
\end{itemize}

\subsubsection{Study note}\label{study-note-1}

Questa formulazione evita di chiudere prematuramente l'obbligatorietà
universale del vocabolario.

\begin{center}\rule{0.5\linewidth}{0.5pt}\end{center}

\subsection{MR-0001 / ADR-0003}\label{mr-0001-adr-0003}

\textbf{Disposition:} Rewritten ThreatForge Decision candidate; proposed
re-parenting

\textbf{Semantic owner:} ThreatForge product capability

\textbf{Proposed parent:} MR-0002 - Sviluppo tracciabile dai requisiti

\textbf{DDTA contract supported / related:} Future governed
realization/verification trace semantics (not yet closed)

\subsubsection{Title}\label{title-2}

Tracciabilità esplicita della realizzazione invece di inferenza dal
repository

\subsubsection{Context}\label{context-2}

ThreatForge deve aiutare a collegare il lavoro tecnico alle conoscenze
governate che lo giustificano. Se la realizzazione viene dedotta
soltanto dalla presenza di file, da naming convention o da prosa
dispersa, il tool non può distinguere in modo affidabile ciò che è
pianificato, ciò che esiste e ciò che costituisce evidenza di verifica.
Una strategia alternativa è mantenere una traccia esplicita e
governabile che il tool possa risolvere senza interpretare
implicitamente il filesystem come fonte semantica.

\subsubsection{Decision}\label{decision-2}

ThreatForge tratta la relazione fra contesto governato, artefatti di
realizzazione ed evidenze di verifica come informazione esplicita
consumata e resa osservabile dal tool. ThreatForge non considera la sola
esistenza di un file né una menzione in prosa come prova sufficiente di
soddisfazione; usa invece le relazioni e le evidenze previste dal
modello/progetto e ne rende interrogabile lo stato.

\subsubsection{Consequences}\label{consequences-2}

\begin{itemize}
\tightlist
\item
  La realizzazione può essere verificata senza ricostruire intenzioni da
  naming o posizione dei file.
\item
  Il tool deve restare allineato alle future semantiche di Requirement,
  realization trace ed evidence invece di inventarle.
\item
  La soluzione richiede dati/relazioni espliciti e quindi più disciplina
  rispetto a inferenze opportunistiche dal repository.
\end{itemize}

\subsubsection{Study note}\label{study-note-2}

La nuova formulazione introduce una vera alternativa di prodotto; non si
limita a dire che ThreatForge `supporta il modello'.

\begin{center}\rule{0.5\linewidth}{0.5pt}\end{center}

\subsection{MR-0001 / ADR-0004}\label{mr-0001-adr-0004}

\textbf{Disposition:} Retained ThreatForge Decision candidate

\textbf{Semantic owner:} ThreatForge tool architecture/support

\textbf{Proposed parent:} MR-0001 - Documentazione di progetto governata
e tracciabile

\textbf{DDTA contract supported / related:} Controlled values owned by
governed model/representation sources

\subsubsection{Title}\label{title-3}

Risoluzione contestuale dei valori controllati da fonti governate

\subsubsection{Context}\label{context-3}

Campi e label controllati possono avere significati diversi in contesti
differenti. Se i consumer ThreatForge codificano enum globali o
interpretano autonomamente stringhe come status, possono attribuire un
significato diverso da quello posseduto dalla fonte governata.
ThreatForge deve scegliere se mantenere interpretazioni locali nei
consumer oppure risolvere i valori nel contesto che li governa.

\subsubsection{Decision}\label{decision-3}

ThreatForge risolve label e valori controllati rispetto alla fonte e al
contesto governato che ne possiedono il significato, invece di mantenere
enum semantici globali nei consumer generici. Il tool può materializzare
cache o proiezioni tecniche, ma tali rappresentazioni devono restare
riconducibili alla fonte contestuale da cui derivano.

\subsubsection{Consequences}\label{consequences-3}

\begin{itemize}
\tightlist
\item
  Consumer differenti possono condividere significati senza duplicare
  inventari locali.
\item
  L'evoluzione dei value set può propagarsi senza reinterpretazioni
  manuali in ogni adapter.
\item
  Resolution e provenance diventano dipendenze esplicite; cache obsolete
  possono applicare significati non correnti.
\end{itemize}

\begin{center}\rule{0.5\linewidth}{0.5pt}\end{center}

\subsection{MR-0001 / ADR-0006}\label{mr-0001-adr-0006}

\textbf{Disposition:} Retained Decision candidate; parent unresolved

\textbf{Semantic owner:} ThreatForge project operations

\textbf{Proposed parent:} Open - do not force into MR-0001 or MR-0002
until MR ownership is retested

\textbf{DDTA contract supported / related:} Recoverability/continuity
evidence; not a universal DDTA contract

\subsubsection{Title}\label{title-4}

Handoff governato per la continuità dello sviluppo ThreatForge

\subsubsection{Context}\label{context-4}

Lo sviluppo di ThreatForge attraversa sessioni di lavoro, programmatori
e LLM differenti. Affidare la continuita alla memoria di una chat o a un
riassunto manuale rende facile perdere Decision, stato del repository,
verifiche e sorgenti necessarie a proseguire il lavoro in modo
riproducibile.

Occorre quindi decidere se la continuita operativa dello sviluppo
ThreatForge resti una pratica informale oppure venga resa ricostruibile
a partire dallo stato governato del progetto.

\subsubsection{Decision}\label{decision-4}

Il processo di sviluppo di ThreatForge usa un \textbf{handoff governato
e ricostruibile} derivato dallo stato canonico del repository e dalle
evidenze operative necessarie a riprendere il lavoro.

L'handoff non sostituisce le fonti governate e non rende canonico il
riassunto della sessione: serve a trasportare abbastanza contesto e
sorgenti per consentire al nuovo agente di verificare e ricostruire lo
stato rilevante.

\subsubsection{Consequences}\label{consequences-4}

\begin{itemize}
\tightlist
\item
  Il cambio di sessione o agente non dipende dalla memoria
  dell'operatore precedente.
\item
  La continuità viene verificata rispetto allo stato del progetto, non
  rispetto a una narrazione isolata.
\item
  L'handoff deve essere mantenuto riproducibile e sufficientemente
  completo.
\item
  Un handoff incompleto o obsoleto può dare falsa confidenza sulla
  continuità.
\end{itemize}

\subsubsection{Study note}\label{study-note-3}

La Decision è significativa; il problema aperto è la collocazione
gerarchica nel set MR ThreatForge.

\begin{center}\rule{0.5\linewidth}{0.5pt}\end{center}

\subsection{MR-0001 / ADR-0007}\label{mr-0001-adr-0007}

\textbf{Disposition:} Retained ThreatForge Decision candidate

\textbf{Semantic owner:} ThreatForge tool architecture

\textbf{Proposed parent:} MR-0001 - Documentazione di progetto governata
e tracciabile

\textbf{DDTA contract supported / related:} L1/L2 governed model and
canonical representation contracts

\subsubsection{Title}\label{title-5}

Contratti eseguibili condivisi per i consumer ThreatForge

\subsubsection{Context}\label{context-5}

Validatori, editor, generatori e migrazioni ThreatForge devono
interpretare lo stesso modello documentale, ma possono farlo in modi
diversi. Se ogni consumer ricodifica autonomamente struttura, valori e
regole, implementazioni formalmente separate possono divergere pur
dichiarando di supportare lo stesso modello. ThreatForge deve quindi
scegliere fra interpretazioni locali per consumer e una rappresentazione
eseguibile comune derivata dalle fonti governate.

\subsubsection{Decision}\label{decision-5}

ThreatForge deriva o materializza contratti eseguibili comuni dai
modelli e dai contratti di rappresentazione governati e fa dipendere da
tali contratti i consumer che necessitano di interpretazione
strutturale. Schemi, validator contract, editor contract e altre viste
tecniche possono differire per formato, ma non diventano inventari
semantici indipendenti mantenuti a mano.

\subsubsection{Consequences}\label{consequences-5}

\begin{itemize}
\tightlist
\item
  Consumer differenti possono condividere struttura e diagnostica senza
  ricodificare il modello in più punti.
\item
  Le proiezioni eseguibili devono restare tracciabili alla fonte
  governata e aggiornabili quando essa cambia.
\item
  Una proiezione obsoleta o arricchita con semantica tool-specifica può
  diventare una fonte concorrente.
\item
  La scelta aumenta la disciplina necessaria per generazione,
  versionamento e verifica delle proiezioni.
\end{itemize}

\subsubsection{Study note}\label{study-note-4}

Questa è una Decision L3 di ThreatForge; il principio generale delle
proiezioni non viene più trattato come invariante L1.

\begin{center}\rule{0.5\linewidth}{0.5pt}\end{center}

\subsection{MR-0001 / ADR-0008}\label{mr-0001-adr-0008}

\textbf{Disposition:} Relocated representation Decision candidate

\textbf{Semantic owner:} DDTA canonical representation contract (L2)

\textbf{Proposed parent:} Not assigned inside the ThreatForge MR
hierarchy

\textbf{DDTA contract supported / related:} Governed entity reference
representation; identity/reference semantics still deferred

\subsubsection{Title}\label{title-6}

Rappresentazione Doc-as-Code canonica dei riferimenti governati

\subsubsection{Context}\label{context-6}

Quando il modello governato permette a un documento di riferirsi a
un'altra entita, la rappresentazione Markdown deve poter distinguere un
riferimento governato da testo ordinario e deve preservare identita
stabile e leggibilita senza ridefinire la semantica dell'entita
referenziata.

Occorre quindi scegliere una rappresentazione Doc-as-Code comune per
serializzare tali riferimenti.

\subsubsection{Decision}\label{decision-6}

La rappresentazione Doc-as-Code studiata per DDTA usa un payload di
riferimento canonico nella forma
\texttt{{[}\textless{}canonical-id\textgreater{}{]}\ \textless{}canonical-title\textgreater{}},
in cui l'identificativo rappresenta l'identita e il titolo e un mirror
leggibile risolto dalla fonte governata.

La rappresentazione di riferimento non crea origine, provenienza,
derivazione o altra relazione semantica: tali relazioni restano definite
dal modello dell'entita e dal contesto che ammette il riferimento.

\subsubsection{Consequences}\label{consequences-6}

\begin{itemize}
\tightlist
\item
  Persone e tool possono leggere e risolvere la stessa rappresentazione
  senza affidarsi a forme libere.
\item
  La serializzazione resta separata da origine, provenienza, derivazione
  e altre relazioni semantiche più forti.
\item
  Parser e authoring devono conoscere le posizioni in cui il payload ha
  significato governato.
\item
  La scelta non viene chiusa definitivamente prima dello studio
  identity/reference/mapping.
\end{itemize}

\begin{center}\rule{0.5\linewidth}{0.5pt}\end{center}

\subsection{MR-0001 / ADR-0009}\label{mr-0001-adr-0009}

\textbf{Disposition:} Removed as standalone Decision candidate

\textbf{Semantic owner:} ThreatForge L4 support/conformance case

\textbf{Proposed parent:} No Decision parent

\textbf{DDTA contract supported / related:} Future Security Requirement
+ Finding contracts

\subsubsection{Disposition rationale}\label{disposition-rationale}

La rewrite non conserva una scelta ThreatForge autonoma sufficiente. Una
volta scelti contratti eseguibili condivisi (ADR-0007) e consumer
estensibili dall'inventario governato (ADR-0010), supportare un futuro
Security Requirement tramite quei contratti è principalmente una prova
di coverage/conformance. La semantica Security Requirement/Finding resta
nelle rispettive fasi L1/L2. Questo slot storico deve quindi diventare
un caso di verifica: ThreatForge dimostra di poter supportare
l'estensione senza hard-code metodologico, parentage o lifecycle non
posseduti dal tool.

\begin{center}\rule{0.5\linewidth}{0.5pt}\end{center}

\subsection{MR-0001 / ADR-0010}\label{mr-0001-adr-0010}

\textbf{Disposition:} Retained ThreatForge Decision candidate

\textbf{Semantic owner:} ThreatForge tool architecture

\textbf{Proposed parent:} MR-0001 - Documentazione di progetto governata
e tracciabile

\textbf{DDTA contract supported / related:} Governed active
document-model inventory/catalog

\subsubsection{Title}\label{title-7}

Inventario dei modelli ThreatForge derivato dal catalogo governato

\subsubsection{Context}\label{context-7}

Quando entrano nuovi tipi documentali, consumer generici possono
scoprire l'inventario attivo tramite liste locali, cardinalità fisse,
filesystem discovery oppure una fonte governata comune. Liste e
cardinalità codificate nei consumer possono divergere dall'inventario
realmente attivo anche se ciascun consumer continua a funzionare.

\subsubsection{Decision}\label{decision-7}

ThreatForge deriva l'inventario dei modelli documentali supportati e dei
relativi contratti dal catalogo governato o da sue proiezioni
deterministiche. I consumer generici non ridefiniscono localmente quali
modelli esistono e non usano cardinalità numeriche come vincolo
architetturale permanente.

\subsubsection{Consequences}\label{consequences-7}

\begin{itemize}
\tightlist
\item
  Nuovi modelli possono entrare attraverso un'estensione esplicita
  invece di una serie di eccezioni scollegate.
\item
  I consumer possono essere verificati contro lo stesso inventario
  attivo.
\item
  L'attivazione richiede coverage dichiarata dei consumer necessari.
\item
  Registrare un modello prima che i consumer richiesti lo supportino può
  produrre falsa estensibilità.
\end{itemize}

\begin{center}\rule{0.5\linewidth}{0.5pt}\end{center}

\subsection{MR-0002 / ADR-0001}\label{mr-0002-adr-0001}

\textbf{Disposition:} Retained ThreatForge Decision candidate

\textbf{Semantic owner:} ThreatForge product architecture

\textbf{Proposed parent:} MR-0002 - Sviluppo tracciabile dai requisiti

\textbf{DDTA contract supported / related:} Multiple DDTA contracts can
be consumed; no single contract defines this architecture

\subsubsection{Title}\label{title-8}

Separazione tra capacità ThreatForge, orchestrazione applicativa e
adapter

\subsubsection{Context}\label{context-8}

ThreatForge deve esporre capacità di supporto alla documentazione e al
lavoro governato attraverso superfici differenti. Se comportamento
applicativo, interpretazione dei contratti e interazione con l'utente
vengono incorporati direttamente in ogni superficie, ciascun adapter può
diventare una variante autonoma del prodotto.

\subsubsection{Decision}\label{decision-8}

ThreatForge separa le capacità riutilizzabili del tool, l'orchestrazione
applicativa e gli adapter di delivery/authoring. Le capacità consumano i
contratti governati necessari; l'applicazione orchestra i casi d'uso;
gli adapter traducono interazione e trasporto senza possedere la logica
applicativa o ridefinire i contratti consumati.

\subsubsection{Consequences}\label{consequences-8}

\begin{itemize}
\tightlist
\item
  CLI, editor, CI e applicazione possono riusare le stesse capacità.
\item
  Cambiare adapter non richiede una nuova interpretazione del modello
  DDTA.
\item
  I confini fra capacità, orchestrazione e delivery devono essere
  mantenuti esplicitamente.
\item
  La scelta del primo adapter concreto resta separata.
\end{itemize}

\subsubsection{Study note}\label{study-note-5}

Rimosso il riferimento normativo a VS Code come prima superficie per non
anticipare ADR-0005.

\begin{center}\rule{0.5\linewidth}{0.5pt}\end{center}

\subsection{MR-0002 / ADR-0002}\label{mr-0002-adr-0002}

\textbf{Disposition:} Retained ThreatForge Decision candidate

\textbf{Semantic owner:} ThreatForge project operations

\textbf{Proposed parent:} MR-0002 - Sviluppo tracciabile dai requisiti

\textbf{DDTA contract supported / related:} Repository-specific
governance, not a universal DDTA contract

\subsubsection{Title}\label{title-9}

Pubblicazione governata del repository ThreatForge

\subsubsection{Context}\label{context-9}

Nel repository di sviluppo ThreatForge, commit e push diretti possono
pubblicare uno stato in cui proiezioni versionate sono obsolete o i gate
locali non sono stati eseguiti. Affidare la sequenza corretta alla
memoria dello sviluppatore rende la governance del repository
facoltativa.

Occorre quindi decidere se la pubblicazione di \textbf{ThreatForge
stesso} resti una sequenza manuale di comandi Git oppure un'operazione
governata del progetto.

\subsubsection{Decision}\label{decision-9}

Il progetto ThreatForge usa una \textbf{operazione governata di
pubblicazione Git} che esegue le precondizioni dichiarate dal repository
prima di stage, commit e push e mantiene distinta una modalita di sola
verifica priva di mutazioni Git.

L'operazione orchestra gate e materializzazioni posseduti dalle
rispettive fonti senza duplicarne la logica.

\subsubsection{Consequences}\label{consequences-9}

\begin{itemize}
\tightlist
\item
  La normale pubblicazione di ThreatForge incorpora le verifiche del
  progetto invece di affidarsi alla disciplina individuale.
\item
  Fonti e proiezioni versionate possono essere pubblicate in uno stato
  coerente.
\item
  Il workflow del repository è più strutturato del Git diretto.
\item
  Un orchestratore errato può introdurre side effect o staging inattesi.
\end{itemize}

\begin{center}\rule{0.5\linewidth}{0.5pt}\end{center}

\subsection{MR-0002 / ADR-0003}\label{mr-0002-adr-0003}

\textbf{Disposition:} Retained ThreatForge Decision candidate

\textbf{Semantic owner:} ThreatForge product capability

\textbf{Proposed parent:} MR-0002 - Sviluppo tracciabile dai requisiti

\textbf{DDTA contract supported / related:} Future realization
trace/lifecycle contracts

\subsubsection{Title}\label{title-10}

Scaffold esplicito come avvio governato della realizzazione

\subsubsection{Context}\label{context-10}

Quando ThreatForge aiuta ad avviare lavoro tecnico da un contesto
governato, la creazione manuale di un file può perdere il collegamento
con ciò che ne giustifica l'esistenza oppure essere interpretata come
realizzazione completata. ThreatForge deve scegliere se limitarsi a
creare artefatti tecnici o rappresentare esplicitamente l'avvio come
operazione distinta dal completamento.

\subsubsection{Decision}\label{decision-10}

ThreatForge offre una operazione di scaffold che crea un punto di
partenza tecnico collegato al contesto governato selezionato e alla
tracciabilità prevista dal progetto. La presenza dello scaffold non
costituisce da sola evidenza di soddisfazione o completamento; eventuali
stati formali e transizioni vengono consumati dai contratti che li
possiedono.

\subsubsection{Consequences}\label{consequences-10}

\begin{itemize}
\tightlist
\item
  L'avvio del lavoro conserva il legame con il contesto governato.
\item
  Uno scaffold non viene confuso con completamento.
\item
  La creazione tecnica richiede una operazione tool-aware aggiuntiva.
\item
  Hard-code di stati o transizioni anticiperebbe semantiche non ancora
  chiuse.
\end{itemize}

\begin{center}\rule{0.5\linewidth}{0.5pt}\end{center}

\subsection{MR-0002 / ADR-0004}\label{mr-0002-adr-0004}

\textbf{Disposition:} Retained ThreatForge Decision candidate

\textbf{Semantic owner:} ThreatForge product capability

\textbf{Proposed parent:} MR-0002 - Sviluppo tracciabile dai requisiti

\textbf{DDTA contract supported / related:} Governed document/model
contracts and future effective-context projections

\subsubsection{Title}\label{title-11}

Authoring ThreatForge guidato dai contratti governati del modello

\subsubsection{Context}\label{context-11}

La sola validazione a posteriori costringe autori e LLM a ricostruire
manualmente identita, relazioni, struttura e valori controllati,
aumentando la probabilita di creare documenti incoerenti prima che i
check li rilevino.

Occorre quindi decidere se ThreatForge limiti il proprio supporto ai
controlli finali oppure usi la conoscenza governata del modello anche
durante l'authoring.

\subsubsection{Decision}\label{decision-11}

ThreatForge offre \textbf{authoring guidato e deterministico} derivando
suggerimenti, generatori e vincoli dai contratti governati del
modello/progetto anziche mantenere copie locali delle regole nei singoli
adapter.

L'authoring riduce gli errori prima della validazione, mentre i check
restano un controllo indipendente finale. Il tool guida l'uso del
modello ma non diventa la fonte della sua semantica.

\subsubsection{Consequences}\label{consequences-11}

\begin{itemize}
\tightlist
\item
  Persone e LLM ricevono contesto rilevante durante la scrittura invece
  di ricordare manualmente tutte le regole.
\item
  CLI/editor possono condividere conoscenza derivata.
\item
  Proiezioni e contratti di authoring devono restare freschi e
  tracciabili alla fonte governata.
\item
  Assistenza obsoleta può guidare verso contenuto incompatibile con il
  modello corrente.
\end{itemize}

\begin{center}\rule{0.5\linewidth}{0.5pt}\end{center}

\subsection{MR-0002 / ADR-0005}\label{mr-0002-adr-0005}

\textbf{Disposition:} Retained ThreatForge Decision candidate

\textbf{Semantic owner:} ThreatForge adapter decision

\textbf{Proposed parent:} MR-0002 - Sviluppo tracciabile dai requisiti

\textbf{DDTA contract supported / related:} ThreatForge capabilities
exposed without duplicating their semantics

\subsubsection{Title}\label{title-12}

VS Code task come adapter locale sottile di ThreatForge

\subsubsection{Context}\label{context-12}

Le capacita ThreatForge disponibili via CLI sono riutilizzabili ma poco
ergonomiche durante il lavoro quotidiano in VS Code. Duplicare regole o
valori nell'editor creerebbe una seconda implementazione, mentre una
integrazione piu complessa non e necessaria per semplici operazioni di
lancio.

Occorre quindi scegliere la prima superficie locale con cui esporre tali
capacita nell'editor.

\subsubsection{Decision}\label{decision-12}

ThreatForge usa un \textbf{catalogo di task VS Code} come adapter locale
sottile sopra gli entrypoint governati esistenti. I task espongono
operazioni e input senza possedere semantica del modello, valori
canonici o logica applicativa e mantengono distinguibili operazioni
read-only e operazioni con side effect.

\subsubsection{Consequences}\label{consequences-12}

\begin{itemize}
\tightlist
\item
  Le operazioni frequenti diventano accessibili dall'editor senza una
  seconda implementazione delle regole.
\item
  La capacità sottostante resta riusabile fuori da VS Code.
\item
  Il catalogo deve seguire l'evoluzione degli entrypoint.
\item
  Regole locali nei task potrebbero trasformare l'adapter in una fonte
  concorrente.
\end{itemize}

\begin{center}\rule{0.5\linewidth}{0.5pt}\end{center}

\subsection{MR-0002 / ADR-0006}\label{mr-0002-adr-0006}

\textbf{Disposition:} Retained ThreatForge Decision candidate

\textbf{Semantic owner:} ThreatForge product architecture

\textbf{Proposed parent:} MR-0002 - Sviluppo tracciabile dai requisiti

\textbf{DDTA contract supported / related:} Canonical document/model
contracts consumed for live assistance

\subsubsection{Title}\label{title-13}

Core ThreatForge condiviso per assistenza Markdown contestuale

\subsubsection{Context}\label{context-13}

I task possono lanciare operazioni ma non interpretano il testo non
salvato e la posizione del cursore. Se VS Code e un editor web
implementassero separatamente completion e diagnostica, ciascun adapter
potrebbe ricostruire una propria versione delle regole documentali.

Occorre quindi decidere dove ThreatForge esegua l'analisi contestuale
usata dall'assistenza live.

\subsubsection{Decision}\label{decision-13}

ThreatForge usa un \textbf{core di assistenza Markdown indipendente
dall'editor e privo di side effect} che consuma i contratti/proiezioni
governate del modello e restituisce informazione di assistenza sul testo
corrente. VS Code e altri editor restano adapter di presentazione e non
possiedono inventari o regole semantiche autonome.

\subsubsection{Consequences}\label{consequences-13}

\begin{itemize}
\tightlist
\item
  Editor differenti possono produrre assistenza semanticamente coerente
  sullo stesso stato documentale.
\item
  Il core può essere verificato indipendentemente dall'interfaccia.
\item
  Il contratto fra core e adapter deve restare stabile.
\item
  Trasformazioni adapter-specifiche possono ancora produrre differenze
  di range, ordine o presentazione.
\end{itemize}

\begin{center}\rule{0.5\linewidth}{0.5pt}\end{center}

\subsection{MR-0002 / ADR-0007}\label{mr-0002-adr-0007}

\textbf{Disposition:} Retained ThreatForge Decision candidate

\textbf{Semantic owner:} ThreatForge product architecture

\textbf{Proposed parent:} MR-0002 - Sviluppo tracciabile dai requisiti

\textbf{DDTA contract supported / related:} No single DDTA semantic
contract; this is internal application architecture

\subsubsection{Title}\label{title-14}

Architettura applicativa ThreatForge indipendente dai delivery adapter

\subsubsection{Context}\label{context-14}

Le capacita di prodotto ThreatForge devono poter essere esposte da CLI e
successivamente da interfacce web. Incorporare comportamento
applicativo, accesso infrastrutturale e traduzione del transport nello
stesso runner renderebbe ogni nuova superficie una migrazione o una
duplicazione.

Occorre quindi scegliere un confine applicativo riutilizzabile interno
al tool.

\subsubsection{Decision}\label{decision-14}

ThreatForge organizza le capacita di prodotto riutilizzabili attorno a
\textbf{servizi applicativi con dipendenze esplicite verso porte},
implementate da adapter infrastrutturali e composte in un confine di
assembly della feature. CLI, HTTP e altre superfici traducono
input/output ma non possiedono il comportamento applicativo.

Questa architettura organizza il tool ThreatForge; non definisce la
struttura interna dei progetti governati ne la semantica del metamodello
DDTA.

\subsubsection{Consequences}\label{consequences-14}

\begin{itemize}
\tightlist
\item
  CLI e web possono condividere comportamento e test applicativi.
\item
  Infrastruttura e transport restano sostituibili dietro confini
  espliciti.
\item
  Una feature richiede responsabilità e assembly più espliciti di un
  runner monolitico.
\item
  Astrazioni premature o leakage di logica negli adapter possono
  annullare il vantaggio.
\end{itemize}

\subsubsection{Study note}\label{study-note-6}

Zod, OpenAPI, React e framework HTTP restano downstream o candidate
Decision separate se acquisiscono trade-off autonomi.

\begin{center}\rule{0.5\linewidth}{0.5pt}\end{center}

\subsection{Corpus-level observations}\label{corpus-level-observations}

\begin{enumerate}
\def\labelenumi{\arabic{enumi}.}
\tightlist
\item
  The corrected set is intentionally asymmetric. No quota of Decisions
  per Macro Requirement is assumed.
\item
  Tool architecture Decisions are allowed to choose a concrete
  realization strategy even when that strategy supports an L1/L2 DDTA
  contract.
\item
  A DDTA contract being supported is not itself sufficient to justify a
  ThreatForge Decision; the ThreatForge document must still contain an
  autonomous choice.
\item
  Historical identifiers are retained only to preserve evidence
  traceability. Relocated candidates will need a new owning hierarchy
  before any canonical migration.
\item
  MR-0001/ADR-0009 demonstrates that some historical ADR content is
  better represented as L4 support/conformance coverage rather than as
  another Decision.
\end{enumerate}

\subsection{Validation target}\label{validation-target}

The next validation step is to apply the Revision 7 Decision tests to
this corrected corpus itself. If a rewritten candidate still lacks one
natural parent, one coherent choice, or non-redundant contribution, the
defect remains in the authoring/corpus rather than being hidden by
historical structure.

\end{document}
